\chapter{Résultats}
\section{Comparaison Avant et Après}
\subsection{Avant (sans pare-feu)}
\begin{lstlisting}
Nmap scan report for 192.168.1.111
Not shown: 1023 closed tcp ports (reset)
PORT STATE SERVICE
22/tcp open ssh
Scan time: 0.22 seconds
\end{lstlisting}
All 1023 blocked ports responded with a TCP RST packet, revealing that no service was listening.

\subsection{Après (avec pare-feu)}
\begin{lstlisting}
Nmap scan report for 192.168.1.111
Not shown: 1021 filtered tcp ports (no-response)
PORT STATE SERVICE
22/tcp open ssh
80/tcp closed http
443/tcp closed https
Scan time: 4.79 seconds
\end{lstlisting}

\section{Analyse des Résultats}
\begin{table}[h!]
\centering
\begin{tabular}{|l|l|l|}
\hline
\textbf{Métrique} & \textbf{Avant} & \textbf{Après} \\
\hline
État des ports bloqués & closed & filtered \\
Durée du scan & 0.22 seconds & 4.79 seconds \\
Ports visibles & 1 (SSH) & 3 (SSH, HTTP, HTTPS) \\
Réponse des sondes & Paquet RST envoyé & Aucune réponse \\
\hline
\end{tabular}
\caption{Comparaison des scans nmap}
\label{tab:nmap_comparison}
\end{table}

L'explication de \texttt{closed} vs \texttt{filtered} :
\begin{itemize}
    \item \texttt{closed} — le serveur a répondu avec un paquet TCP RST.
    \item \texttt{filtered} — le pare-feu a silencieusement rejeté le paquet. Cela a forcé nmap à attendre des délais d'attente, augmentant considérablement la durée du scan.
\end{itemize}
\lipsum[8]
